\section{常微分方程}

\subsection{线性微分方程和非线性微分方程}

\begin{definition}[线性微分方程]
    对于一个微分方程 $F(y, y', y'', \dots) + C = 0$，若 $F$ 关于 $y, y', y'', \dots$ 都是线性函数，那么称该微分方程为\textbf{线性微分方程}；否则称该微分方程为\textbf{非线性微分方程}。

    特别地，当 $C = 0$ 时，我们称该微分方程为\textbf{齐次线性微分方程}。
\end{definition}

微分方程中含有的最高阶导数的阶数称为微分方程的\textbf{阶}。

\subsection{微分方程的解}

代入微分方程能够使方程成立的函数 $\varphi(x)$ 称为微分方程的\textbf{解}。

若微分方程的解中含有的任意常数的个数与方程阶数相同，且这些常数之间不能合并，那么称这个解为该微分方程的\textbf{通解}。

用于确定通解中任意常数的值的条件称为\textbf{初始条件}，当通解中所有常数的值确定后，我们称这个解为该微分方程的\textbf{特解}。

\subsection{分离变量方程}

\begin{definition}[分离变量方程]
    对于形式为 $y' = f(x) g(y)$ 的微分方程，称该微分方程为\textbf{分离变量方程}。
\end{definition}

对于这类方程，我们可以写成一个更简洁的形式：

\begin{equation}
    g(y) \dd{y} = f(x) \dd{x} \tag{1}
\end{equation}


设 $f(x), g(x)$ 是连续的，$y = \varphi(x)$ 是方程 $(1)$ 的解，那么
$$
g(\varphi(x)) \varphi'(x) = f(x)
$$
两端同时积分，得到：
$$
\int g(\varphi(x)) \varphi'(x) \dd{x} = \int f(x) \dd{x}
$$
进行换元，于是有：
$$
\int g(y) \dd{y} = \int f(x) \dd{x}
$$
设 $f(x), g(y)$ 的原函数分别为 $F(x), G(y)$，那么有：

\begin{equation}
    G(y) = F(x) + C \tag{2}
\end{equation}

根据方程 $(2)$ 即可求得方程 $(1)$ 的通解。
